\chapter{Celui qu'on ne surprenait jamais}
\label{chap:celui-quon-ne-surprenait-jamais}

Les mois continuaient de passer et, parmi toutes les choses que Kael apprenait sur les
routes, certaines l'intéressaient plus que d'autres.

Il avait toujours su qu'il était possible de se déplacer discrètement. Au
village, les chasseurs lui avaient appris à éviter les branches mortes, à
poser les pieds avec précaution et à approcher certaines proies sans les faire
fuir. Mais il s'agissait surtout de chasse. On cherchait à ne pas être entendu
par un animal qui ne savait pas que l'on était là.

Avec Caliban Dalsen, le problème était très différent.

Cal savait que Kael était là.

Même lorsqu'il n'aurait pas dû.

\scenebreak

Kael le remarqua quelques jours après son arrivée dans la caravane.

Cal était occupé près d'un chariot, le dos tourné, lorsqu'il décida de
s'approcher sans se faire entendre.

Il n'avait pas vraiment de raison de le faire.

Simplement l'envie de voir s'il en était capable.

Kael ralentit donc ses pas, contourna quelques caisses et prit soin d'éviter
les endroits où le sol risquait de craquer sous son poids.

Il n'était plus qu'à quelques pas lorsque Cal parla sans même se retourner.

— Tu veux quelque chose ?

Kael s'arrêta.

Cal continua tranquillement ce qu'il était en train de faire.

— Non.

— Alors pourquoi tu me suis ?

Kael fronça les sourcils.

— Je ne te suivais pas.

Cette fois, Cal se retourna.

Il regarda Kael.

Puis l'endroit d'où celui-ci venait.

Puis de nouveau Kael.

— Bien sûr.

Kael repartit sans rien ajouter.

Il recommença le lendemain.

Cette fois, Cal tourna la tête avant même qu'il soit arrivé à proximité.

Kael essaya encore deux jours plus tard.

Même résultat.

Puis une quatrième fois.

Cal leva simplement une main dans sa direction sans interrompre la
conversation qu'il était en train d'avoir avec un autre voyageur.

Cela commençait à devenir agaçant.

Évidemment, plus Kael échouait, plus il voulait réussir.

\scenebreak

Il commença à préparer ses tentatives.

Approcher Cal directement ne fonctionnait pas.

Il essaya donc de profiter des chariots pour disparaître de son champ de
vision.

Cal le vit quand même.

Il attendit qu'un groupe de voyageurs passe entre eux avant de changer de
direction.

Cal se retourna.

Il tenta de se rapprocher pendant que plusieurs personnes déchargeaient des
caisses et profita du bruit pour masquer ses pas.

Cal leva les yeux exactement au moment où il arrivait.

Kamatsu finit par remarquer son manège.

— Qu'est-ce que tu fais ?

Kael, accroupi derrière un chariot, lui fit immédiatement signe de se taire.

Kamatsu regarda dans la direction qu'il observait.

Cal était assis un peu plus loin.

Il regarda Kael.

Puis Cal.

Puis de nouveau Kael.

— Tu espionnes Cal ?

— Chut.

Kamatsu baissa la voix.

— Pourquoi ?

— J'essaie de le surprendre.

— Pourquoi ?

Kael réfléchit quelques secondes.

Il n'avait pas réellement de réponse à cette question.

— Parce que j'y arrive pas.

Kamatsu sembla considérer cette explication.

— D'accord.

Il s'assit à côté de lui.

Kael tourna la tête.

— Qu'est-ce que tu fais ?

— Je regarde.

— Tu vas me faire repérer.

— Je ne fais rien.

— Justement.

Kamatsu sourit.

— Vas-y.

Kael attendit encore quelques secondes avant de quitter sa cachette.

Il contourna silencieusement le chariot, profita du bruit d'une caisse que l'on
déplaçait un peu plus loin et se glissa derrière une pile de marchandises.

Cal n'avait pas bougé.

Kael continua.

Quelques pas encore.

Cette fois, il allait réussir.

Il contourna le dernier chariot.

Cal n'était plus là.

Kael s'immobilisa.

Il regarda à gauche.

Puis à droite.

— Tu cherches quelqu'un ?

Kael sursauta.

Cal se tenait derrière lui.

Un rire éclata près du chariot.

Kael tourna la tête.

Kamatsu avait manifestement beaucoup moins de difficultés à apprécier la
situation que lui.

Cal regarda dans sa direction.

— Ton frère est plus discret que toi.

Kamatsu cessa immédiatement de rire.

Kael sourit.

Cela compensait un peu.

\scenebreak

Après quelques secondes passées à essayer de comprendre comment Cal avait pu
se retrouver derrière lui, Kael abandonna.

Mais cette fois, au lieu de simplement préparer une nouvelle tentative, il
alla le retrouver un peu plus tard.

— Caliban Dalsen ?

L'homme leva les yeux.

— Cal.

Kael attendit.

— Quoi ?

— Appelle-moi Cal.

— Pourquoi ?

— Parce que si tu cries « Caliban Dalsen » chaque fois que tu veux me parler,
toute la caravane saura où je suis.

Kael dut reconnaître que l'argument était difficilement contestable.

Il s'assit près de lui.

— Comment tu savais que je te suivais ?

Cal sourit.

— Laquelle des fois ?

Kael fronça les sourcils.

— Aujourd'hui.

— Ton frère te regardait depuis l'autre côté du camp.

Kael tourna immédiatement la tête vers Kamatsu.

— Mais avant ?

— Avant quoi ?

— Les autres fois.

Cal prit quelques secondes avant de répondre.

— Parce que tu fais beaucoup trop d'efforts pour ne pas faire de bruit.

Kael le regarda sans comprendre.

— Je ne faisais presque aucun bruit.

— Justement.

Cal se releva.

— Viens.

\scenebreak

Il conduisit Kael jusqu'à l'extérieur du camp.

— Marche.

Kael obéit.

Il posa soigneusement le pied devant lui.

Puis l'autre.

Il évita une branche.

Contourna quelques feuilles sèches.

Cal l'arrêta.

— Tu fais quoi ?

— Je marche sans faire de bruit.

— Non. Tu essayes de marcher sans faire de bruit.

Kael attendit la suite.

Cal désigna les environs.

Le vent faisait bouger les branches au-dessus d'eux. Des oiseaux chantaient
plus loin. Derrière eux, la caravane produisait son mélange habituel de voix,
de roues, d'animaux et d'objets déplacés.

— Le silence n'existe pas, expliqua Cal. Alors arrête de le chercher.

Il ramassa une petite branche et la brisa entre ses doigts.

Kael regarda autour de lui.

— Tout fait du bruit.

— Exactement.

Cal commença à marcher.

Il n'était pas silencieux.

C'était précisément ce qui surprit Kael.

Ses pas produisaient parfois un léger froissement. Une feuille bougeait sous
sa botte. Une branche se déplaçait lorsqu'il passait.

Mais aucun de ces sons ne semblait inhabituel.

Ils disparaissaient simplement parmi les autres.

— Si tu attends que tout soit parfaitement silencieux pour avancer, poursuivit
Cal, le moindre bruit que tu feras sera le seul qu'on entendra.

Il attendit qu'une rafale traverse les arbres, puis posa le pied sur un tapis
de feuilles.

Le froissement disparut dans celui des branches.

Kael comprit.

Ce fut probablement la première véritable leçon de Cal.

\scenebreak

Cal ne devint jamais réellement son maître.

Il n'organisait pas de leçons comme Robin avait pu le faire avec son arc et ne
semblait pas particulièrement intéressé par l'idée d'avoir un élève.

Mais lorsqu'il voyait Kael essayer quelque chose, il lui arrivait de le
corriger.

Il lui montra où poser les pieds lorsque le sol était irrégulier.

Comment répartir son poids avant de déplacer l'autre jambe.

Comment utiliser les bruits environnants plutôt que de chercher à les éviter.

Comment suivre quelqu'un sans marcher exactement dans ses traces.

Mais surtout, il lui apprit que la discrétion ne concernait pas uniquement les
pas.

Un après-midi, Cal traversait un marché avec la caravane lorsqu'il s'arrêta
brusquement devant un étal.

Kael, qui le suivait depuis plusieurs minutes, s'arrêta lui aussi.

Cal se retourna.

— Perdu.

Kael soupira.

— Je n'ai pas fait de bruit.

— Je sais.

— Alors comment tu m'as vu ?

— Je ne t'ai pas vu.

Kael attendit.

Cal désigna une échoppe quelques mètres plus loin.

— Tu regardes quoi ?

— Toi.

— Voilà.

Kael fronça les sourcils.

— Si tu passes dix minutes à regarder quelqu'un, poursuivit Cal, il finira par
le sentir. Et même s'il ne te voit pas, quelqu'un d'autre remarquera peut-être
que tu passes ton temps à le fixer.

— Alors je regarde où ?

— Partout.

Cal reprit sa marche.

— Sauf là où tu veux regarder.

Kael resta quelques secondes immobile avant de le suivre.

Cette leçon-là lui demanda beaucoup plus de temps.

\scenebreak

Les semaines suivantes, Kael continua à essayer de surprendre Cal.

Il changea ses méthodes.

Il cessa de chercher le silence absolu.

Apprit à attendre.

À écouter.

À utiliser le passage d'un chariot, une conversation trop forte ou une rafale
de vent.

Il apprit également à regarder sans fixer et à choisir ses déplacements en
fonction de ce que faisaient les autres plutôt que seulement en fonction de
l'endroit où il voulait aller.

Ses progrès étaient réels.

Cela ne changea absolument rien au résultat.

Un soir, il parvint à s'approcher suffisamment près de Cal pour être certain
d'avoir enfin réussi.

Il tendit lentement la main vers son épaule.

— Non.

Kael s'immobilisa.

Cal ne s'était même pas retourné.

— Comment ?

— Trop facile.

— Mais comment ?

Cal tourna enfin la tête.

Il souriait.

— Une autre fois.

Kael soupira.

Il n'y arriva jamais.

Pas une seule fois.

Mais avec le temps, quelque chose changea malgré tout.

Au début, lorsqu'il échouait, Kael se demandait uniquement comment Cal avait
pu savoir qu'il était là.

Puis il commença à comprendre ses erreurs avant même que Cal les lui explique.

Un pas posé trop vite.

Un mouvement effectué au mauvais moment.

Un regard maintenu trop longtemps.

Une impatience qui l'avait poussé à avancer alors qu'il aurait simplement dû
attendre.

Kael ne parvenait toujours pas à surprendre Cal.

Mais il commençait enfin à comprendre pourquoi.